% !TEX TS-program = xelatex
% !TEX encoding = UTF-8 Unicode
% !Mode:: "TeX:UTF-8"

\documentclass{resume}

\begin{document}
\pagenumbering{gobble}

\name{你的姓名}

\basicInfo{
  \email{you@example.com} \textperiodcentered\
  \phone{(+00) 0000000000} \textperiodcentered\
  \github{github.com/your-id}
}

\headline{目标岗位 \textperiodcentered\ N 年工作经验 \textperiodcentered\ 方向 A / 方向 B / 工程化落地}

\section{\faUser\ 个人概览}
\begin{itemize}
  \item 具备面向业务系统、数据驱动决策和生产环境交付的实践经验，能够在目标、资源和工程约束之间做取舍。
  \item 擅长将模糊业务问题拆解为技术方案、实验指标、可复用系统和团队实践，关注长期可维护性与持续迭代效率。
\end{itemize}

\section{\faGraduationCap\ 教育背景}

\datedline{\makebox[6.2em][l]{\textbf{某某大学}}\makebox[5.2em][l]{\textit{硕士}}计算机科学与技术}{\textbf{YYYY.MM -- YYYY.MM}}
\datedline{\makebox[6.2em][l]{\textbf{某某大学}}\makebox[5.2em][l]{\textit{本科}}软件工程}{\textbf{YYYY.MM -- YYYY.MM}}

\section{\faBriefcase\ 工作经历}

\datedsubsection{\textbf{公司 A} \quad {\normalfont 目标岗位 / 当前职级}}{YYYY.MM -- 至今}
\begin{itemize}
  \item 面向高流量业务场景，负责核心系统、算法策略或平台能力建设，平衡产品目标、数据约束、工程稳定性与线上效果。
  \item 搭建端到端迭代流程，覆盖数据构造、离线评估、实验分析、上线发布、监控告警与效果复盘，提升团队持续迭代效率。
  \item 将一次性规则、脚本或人工流程沉淀为可复用模块、服务接口和评估标准，减少重复工作并提升问题定位效率。
  \item 与产品、工程、数据、运营等角色协作，拆解模糊需求，推动方案从原型验证进入稳定生产系统。
\end{itemize}

\datedsubsection{\textbf{公司 B} \quad {\normalfont 岗位名称}}{YYYY.MM -- YYYY.MM}
\begin{itemize}
  \item 独立负责跨团队项目从问题定义到上线落地，包括需求分析、方案设计、原型验证、实验评估与发布计划。
  \item 建设通用组件、分析看板、调试工具或数据链路，降低后续项目接入成本，并提升实验 review 与问题排查效率。
  \item 结合线上指标、用户反馈和业务约束定位问题，将分析结论转化为具体技术改动，并沉淀为团队文档或最佳实践。
\end{itemize}

\datedsubsection{\textbf{公司 C} \quad {\normalfont 早期岗位名称}}{YYYY.MM -- YYYY.MM}
\begin{itemize}
  \item 参与后端、数据、平台、模型或客户端功能研发，关注功能正确性、代码可维护性、交付节奏和线上稳定性。
  \item 早期经历建议保持精简：只保留能支撑当前目标岗位的代表性项目，弱相关细节可以压缩或删除。
\end{itemize}

\section{\faFileTextO\ 项目 / 论文 / 专利}
\begin{itemize}
  \item 代表项目：构建内部工具、开源项目、技术平台或应用研究项目，用一句话说明问题背景、技术方案和业务价值。
  \item 论文或专利：只保留和目标岗位相关、能证明专业深度的内容；如有公开链接，可以附 DOI、专利号或项目地址。
  \item 技术影响力：可放精选技术文章、公开演讲、内部分享或方法论沉淀，突出实现之外的判断力和表达能力。
\end{itemize}

\section{\faCogs\ 专业能力}
\begin{itemize}
  \item 业务与算法：问题建模、策略优化、搜索 / 推荐 / 排序、实验评估、数据分析、产品迭代与效果复盘。
  \item 工程能力：Python、C++、SQL、Linux、服务端开发、分布式系统、可观测性、问题排查、性能与稳定性优化。
  \item 协作交付：需求拆解、技术规划、实验 review、事故复盘、跨团队推进、文档沉淀、团队 mentoring。
\end{itemize}

\end{document}
